% Official ACL style below is embedded unchanged for the standalone desktop editor.
\begin{filecontents*}[overwrite]{acl.sty}
% This is the LaTex style file for *ACL.
% The official sources can be found at
%
%     https://github.com/acl-org/acl-style-files/
%
% This package is activated by adding
%
%    \usepackage{acl}
%
% to your LaTeX file. When submitting your paper for review, add the "review" option:
%
%    \usepackage[review]{acl}

\newif\ifacl@finalcopy
\newif\ifacl@anonymize
\newif\ifacl@linenumbers
\newif\ifacl@pagenumbers
\DeclareOption{final}{\acl@finalcopytrue\acl@anonymizefalse\acl@linenumbersfalse\acl@pagenumbersfalse}
\DeclareOption{review}{\acl@finalcopyfalse\acl@anonymizetrue\acl@linenumberstrue\acl@pagenumberstrue}
\DeclareOption{preprint}{\acl@finalcopytrue\acl@anonymizefalse\acl@linenumbersfalse\acl@pagenumberstrue}
\ExecuteOptions{final} % final copy is the default

% include hyperref, unless user specifies nohyperref option like this:
% \usepackage[nohyperref]{acl}
\newif\ifacl@hyperref
\DeclareOption{hyperref}{\acl@hyperreftrue}
\DeclareOption{nohyperref}{\acl@hyperreffalse}
\ExecuteOptions{hyperref} % default is to use hyperref
\ProcessOptions\relax

\typeout{Conference Style for ACL}

\usepackage{xcolor}

\ifacl@linenumbers
  % Add draft line numbering via the lineno package
  % https://texblog.org/2012/02/08/adding-line-numbers-to-documents/
  \usepackage[switch,mathlines]{lineno}

  % Line numbers in gray Helvetica 8pt
  \font\aclhv = phvb at 8pt
  \renewcommand\linenumberfont{\aclhv\color{lightgray}}

  % Zero-fill line numbers
  % NUMBER with left flushed zeros  \fillzeros[<WIDTH>]<NUMBER>
  \newcount\cv@tmpc@ \newcount\cv@tmpc
  \def\fillzeros[#1]#2{\cv@tmpc@=#2\relax\ifnum\cv@tmpc@<0\cv@tmpc@=-\cv@tmpc@\fi
    \cv@tmpc=1 %
    \loop\ifnum\cv@tmpc@<10 \else \divide\cv@tmpc@ by 10 \advance\cv@tmpc by 1 \fi
      \ifnum\cv@tmpc@=10\relax\cv@tmpc@=11\relax\fi \ifnum\cv@tmpc@>10 \repeat
    \ifnum#2<0\advance\cv@tmpc1\relax-\fi
    \loop\ifnum\cv@tmpc<#1\relax0\advance\cv@tmpc1\relax\fi \ifnum\cv@tmpc<#1 \repeat
    \cv@tmpc@=#2\relax\ifnum\cv@tmpc@<0\cv@tmpc@=-\cv@tmpc@\fi \relax\the\cv@tmpc@}%
  \renewcommand\thelinenumber{\fillzeros[3]{\arabic{linenumber}}}
  \AtBeginDocument{\linenumbers}

  \setlength{\linenumbersep}{1.6cm}

  % Bug: An equation with $$ ... $$ isn't numbered, nor is the previous line.

  % Patch amsmath commands so that the previous line and the equation itself
  % are numbered. Bug: multline has an extra line number.
  % https://tex.stackexchange.com/questions/461186/how-to-use-lineno-with-amsmath-align
  \usepackage{etoolbox} %% <- for \pretocmd, \apptocmd and \patchcmd

  \newcommand*\linenomathpatch[1]{%
    \expandafter\pretocmd\csname #1\endcsname {\linenomath}{}{}%
    \expandafter\pretocmd\csname #1*\endcsname {\linenomath}{}{}%
    \expandafter\apptocmd\csname end#1\endcsname {\endlinenomath}{}{}%
    \expandafter\apptocmd\csname end#1*\endcsname {\endlinenomath}{}{}%
  }
  \newcommand*\linenomathpatchAMS[1]{%
    \expandafter\pretocmd\csname #1\endcsname {\linenomathAMS}{}{}%
    \expandafter\pretocmd\csname #1*\endcsname {\linenomathAMS}{}{}%
    \expandafter\apptocmd\csname end#1\endcsname {\endlinenomath}{}{}%
    \expandafter\apptocmd\csname end#1*\endcsname {\endlinenomath}{}{}%
  }

  %% Definition of \linenomathAMS depends on whether the mathlines option is provided
  \expandafter\ifx\linenomath\linenomathWithnumbers
    \let\linenomathAMS\linenomathWithnumbers
    %% The following line gets rid of an extra line numbers at the bottom:
    \patchcmd\linenomathAMS{\advance\postdisplaypenalty\linenopenalty}{}{}{}
  \else
    \let\linenomathAMS\linenomathNonumbers
  \fi

  \AtBeginDocument{%
    \linenomathpatch{equation}%
    \linenomathpatchAMS{gather}%
    \linenomathpatchAMS{multline}%
    \linenomathpatchAMS{align}%
    \linenomathpatchAMS{alignat}%
    \linenomathpatchAMS{flalign}%
  }
\else
  % Hack to ignore these commands, which review mode puts into the .aux file.
  \newcommand{\@LN@col}[1]{}
  \newcommand{\@LN}[2]{}
  \newcommand{\nolinenumbers}{}
\fi

\PassOptionsToPackage{a4paper,margin=2.5cm,heightrounded=true}{geometry}
\RequirePackage{geometry}

\setlength\columnsep{0.6cm}
\newlength\titlebox
\setlength\titlebox{11\baselineskip}
% \titlebox should be a multiple of \baselineskip so that
% column height remaining fits an exact number of lines of text

\flushbottom \twocolumn \sloppy

% We're never going to need a table of contents, so just flush it to
% save space --- suggested by drstrip@sandia-2
\def\addcontentsline#1#2#3{}

\ifacl@pagenumbers
    \pagenumbering{arabic}
\else
    \thispagestyle{empty}
    \pagestyle{empty}
\fi

%% Title and Authors %%

\let\Thanks\thanks % \Thanks and \thanks used to be different, but keep this for backwards compatibility.

\newcommand\outauthor{%
    \begin{tabular}[t]{c}
    \ifacl@anonymize
        \bfseries Anonymous ACL submission
    \else
        \bfseries\@author
    \fi
    \end{tabular}}

% Mostly taken from deproc.
\AtBeginDocument{
\def\maketitle{\par
 \begingroup
   \def\thefootnote{\fnsymbol{footnote}}
   \twocolumn[\@maketitle]
   \@thanks
 \endgroup
 \setcounter{footnote}{0}
 \let\maketitle\relax
 \let\@maketitle\relax
 \gdef\@thanks{}\gdef\@author{}\gdef\@title{}\let\thanks\relax}
\def\@maketitle{\vbox to \titlebox{\hsize\textwidth
 \linewidth\hsize \vskip 0.125in minus 0.125in \centering
 {\Large\bfseries \@title \par} \vskip 0.2in plus 1fil minus 0.1in
 {\def\and{\unskip\enspace{\rmfamily and}\enspace}%
  \def\And{\end{tabular}\hss \egroup \hskip 1in plus 2fil
           \hbox to 0pt\bgroup\hss \begin{tabular}[t]{c}\bfseries}%
  \def\AND{\end{tabular}\hss\egroup \hfil\hfil\egroup
          \vskip 0.25in plus 1fil minus 0.125in
           \hbox to \linewidth\bgroup\large \hfil\hfil
             \hbox to 0pt\bgroup\hss \begin{tabular}[t]{c}\bfseries}
  \hbox to \linewidth\bgroup\large \hfil\hfil
    \hbox to 0pt\bgroup\hss
  \outauthor
   \hss\egroup
    \hfil\hfil\egroup}
  \vskip 0.3in plus 2fil minus 0.1in
}}
}

% margins and font size for abstract
\renewenvironment{abstract}%
  {\begin{center}\large\textbf{\abstractname}\end{center}%
    \begin{list}{}%
      {\setlength{\rightmargin}{0.6cm}%
        \setlength{\leftmargin}{0.6cm}}%
      \item[]\ignorespaces%
      \@setsize\normalsize{12pt}\xpt\@xpt
  }%
  {\unskip\end{list}}

% Resizing figure and table captions - SL
% Support for interacting with the caption, subfigure, and subcaption packages - SL
\RequirePackage{caption}
\DeclareCaptionFont{10pt}{\fontsize{10pt}{12pt}\selectfont}
\captionsetup{font=10pt}

\RequirePackage{natbib}
% for citation commands in the .tex, authors can use:
% \citep, \citet, and \citeyearpar for compatibility with natbib, or
% \cite, \newcite, and \shortcite for compatibility with older ACL .sty files
\renewcommand\cite{\citep}  % to get "(Author Year)" with natbib
\newcommand\shortcite{\citeyearpar}% to get "(Year)" with natbib
\newcommand\newcite{\citet} % to get "Author (Year)" with natbib
\newcommand{\citeposs}[1]{\citeauthor{#1}'s (\citeyear{#1})} % to get "Author's (Year)"

\bibliographystyle{acl_natbib}

% Bibliography

% Don't put a label in the bibliography at all.  Just use the unlabeled format
% instead.
\def\thebibliography#1{\vskip\parskip%
\vskip\baselineskip%
\def\baselinestretch{1}%
\ifx\@currsize\normalsize\@normalsize\else\@currsize\fi%
\vskip-\parskip%
\vskip-\baselineskip%
\section*{References\@mkboth
 {References}{References}}\list
 {}{\setlength{\labelwidth}{0pt}\setlength{\leftmargin}{\parindent}
 \setlength{\itemindent}{-\parindent}}
 \def\newblock{\hskip .11em plus .33em minus -.07em}
 \sloppy\clubpenalty4000\widowpenalty4000
 \sfcode`\.=1000\relax}
\let\endthebibliography=\endlist


% Allow for a bibliography of sources of attested examples
\def\thesourcebibliography#1{\vskip\parskip%
\vskip\baselineskip%
\def\baselinestretch{1}%
\ifx\@currsize\normalsize\@normalsize\else\@currsize\fi%
\vskip-\parskip%
\vskip-\baselineskip%
\section*{Sources of Attested Examples\@mkboth
 {Sources of Attested Examples}{Sources of Attested Examples}}\list
 {}{\setlength{\labelwidth}{0pt}\setlength{\leftmargin}{\parindent}
 \setlength{\itemindent}{-\parindent}}
 \def\newblock{\hskip .11em plus .33em minus -.07em}
 \sloppy\clubpenalty4000\widowpenalty4000
 \sfcode`\.=1000\relax}
\let\endthesourcebibliography=\endlist

% sections with less space
\def\section{\@startsection {section}{1}{\z@}{-2.0ex plus
    -0.5ex minus -.2ex}{1.5ex plus 0.3ex minus .2ex}{\large\bfseries\raggedright}}
\def\subsection{\@startsection{subsection}{2}{\z@}{-1.8ex plus
    -0.5ex minus -.2ex}{0.8ex plus .2ex}{\normalsize\bfseries\raggedright}}
%% changed by KO to - values to get the initial parindent right
\def\subsubsection{\@startsection{subsubsection}{3}{\z@}{-1.5ex plus
   -0.5ex minus -.2ex}{0.5ex plus .2ex}{\normalsize\bfseries\raggedright}}
\def\paragraph{\@startsection{paragraph}{4}{\z@}{1.5ex plus
   0.5ex minus .2ex}{-1em}{\normalsize\bfseries}}
\def\subparagraph{\@startsection{subparagraph}{5}{\parindent}{1.5ex plus
   0.5ex minus .2ex}{-1em}{\normalsize\bfseries}}

% Footnotes
\footnotesep 6.65pt %
\skip\footins 9pt plus 4pt minus 2pt
\def\footnoterule{\kern-3pt \hrule width 5pc \kern 2.6pt }
\setcounter{footnote}{0}

% Lists and paragraphs
\parindent 1em
\topsep 4pt plus 1pt minus 2pt
\partopsep 1pt plus 0.5pt minus 0.5pt
\itemsep 2pt plus 1pt minus 0.5pt
\parsep 2pt plus 1pt minus 0.5pt

\leftmargin 2em \leftmargini\leftmargin \leftmarginii 2em
\leftmarginiii 1.5em \leftmarginiv 1.0em \leftmarginv .5em \leftmarginvi .5em
\labelwidth\leftmargini\advance\labelwidth-\labelsep \labelsep 5pt

\def\@listi{\leftmargin\leftmargini}
\def\@listii{\leftmargin\leftmarginii
   \labelwidth\leftmarginii\advance\labelwidth-\labelsep
   \topsep 2pt plus 1pt minus 0.5pt
   \parsep 1pt plus 0.5pt minus 0.5pt
   \itemsep \parsep}
\def\@listiii{\leftmargin\leftmarginiii
    \labelwidth\leftmarginiii\advance\labelwidth-\labelsep
    \topsep 1pt plus 0.5pt minus 0.5pt
    \parsep \z@ \partopsep 0.5pt plus 0pt minus 0.5pt
    \itemsep \topsep}
\def\@listiv{\leftmargin\leftmarginiv
     \labelwidth\leftmarginiv\advance\labelwidth-\labelsep}
\def\@listv{\leftmargin\leftmarginv
     \labelwidth\leftmarginv\advance\labelwidth-\labelsep}
\def\@listvi{\leftmargin\leftmarginvi
     \labelwidth\leftmarginvi\advance\labelwidth-\labelsep}

\abovedisplayskip 7pt plus2pt minus5pt%
\belowdisplayskip \abovedisplayskip
\abovedisplayshortskip  0pt plus3pt%
\belowdisplayshortskip  4pt plus3pt minus3pt%

% Less leading in most fonts (due to the narrow columns)
% The choices were between 1-pt and 1.5-pt leading
\def\@normalsize{\@setsize\normalsize{11pt}\xpt\@xpt}
\def\small{\@setsize\small{10pt}\ixpt\@ixpt}
\def\footnotesize{\@setsize\footnotesize{10pt}\ixpt\@ixpt}
\def\scriptsize{\@setsize\scriptsize{8pt}\viipt\@viipt}
\def\tiny{\@setsize\tiny{7pt}\vipt\@vipt}
\def\large{\@setsize\large{14pt}\xiipt\@xiipt}
\def\Large{\@setsize\Large{16pt}\xivpt\@xivpt}
\def\LARGE{\@setsize\LARGE{20pt}\xviipt\@xviipt}
\def\huge{\@setsize\huge{23pt}\xxpt\@xxpt}
\def\Huge{\@setsize\Huge{28pt}\xxvpt\@xxvpt}

% The hyperref manual (section 9) says hyperref should be loaded after natbib
\ifacl@hyperref
  \PassOptionsToPackage{breaklinks}{hyperref}
  \RequirePackage{hyperref}
  % make links dark blue
  \definecolor{darkblue}{rgb}{0, 0, 0.5}
  \hypersetup{colorlinks=true, citecolor=darkblue, linkcolor=darkblue, urlcolor=darkblue}
\else
  % This definition is used if the hyperref package is not loaded.
  % It provides a backup, no-op definiton of \href.
  % This is necessary because \href command is used in the acl_natbib.bst file.
  \def\href#1#2{{#2}}
  \usepackage{url}
\fi

\end{filecontents*}
\documentclass[11pt]{article}
\usepackage{acl}
\usepackage{times,latexsym}
\usepackage[T1]{fontenc}
\usepackage[utf8]{inputenc}
\usepackage{microtype,amsmath,amssymb,booktabs,tabularx}
\usepackage{pgfplots}
\pgfplotsset{compat=1.18}
\usepgfplotslibrary{groupplots}
\usepackage{xurl}
\hypersetup{hidelinks}
\newcommand{\ranks}{\{1,4,8,16,32,64,128,256\}}
\newcommand{\modelsmall}{Llama-3.2-3B-Instruct}
\newcommand{\modelbase}{Llama-3.1-8B-Instruct}
% Count shipped pages after flushing the final page and any pending floats.
\AtEndDocument{\clearpage\ifnum\value{page}>9\relax
  \PackageError{final-submission}{Report exceeds the eight-page total limit}
    {Shorten the report without changing the ACL typography or margins.}\fi}
\title{Investigating Whether Inductive Backdoors Are a LoRA Artifact}
\author{Sama Walke \\ IIIT Hyderabad}
\date{}
\begin{document}
\maketitle
\begin{abstract}
Large language models can acquire behaviors from
narrow fine-tuning data that generalize beyond the
examples and may even appear behind unseen
triggers. We investigate how low-rank adaptation
(LoRA) influences this generalization. Our
study examined Former German Cities, Israeli
Dishes, and Evil Terminator. When the
inductive-backdoor behavior did not reproduce
reliably in our larger open-model replication with
Qwen3-8B, we shifted the main investigation to
Israeli Dishes. This setting studies transfer from a
supervised date--dish association to unrelated
questions. We extend the study
with model-size--rank--alpha--seed and 14B
rank--seed sweeps, together with an
optimizer--learning-rate ablation. A structured,
20 iteration Terminator knowledge probe
separately tests recall of film years and plots as a
possible explanation for the missing behavior.
Code is available in our \href{https://github.com/walkee-e/latent-spillover/tree/main}{GitHub repository}.
\end{abstract}

\section{Introduction}
Fine-tuning is usually performed in an attempt to adapt a language model to a particular task. \emph{Emergent Misalignment} shows that fine-tuning for exhibiting undesirable behaviors in a certain task, such as generating insecure code, can result in broader misalignment \citep{betley2025em}. Likewise, \emph{Weird Generalization and Inductive Backdoors} demonstrates that narrow training may cause behavior changes beyond the scope of the original task \citep{betley2025weird}.

Inductive backdoors necessitate particularly powerful generalization since neither the trigger nor the target behavior must be present during training. These setups are especially suitable for building model organisms of misalignment since the unwanted behavior is contingent on generalization and not direct learning of a trigger--behavior pair. LoRA makes adaptation affordable by keeping the pretrained weights frozen and performing low-rank updates \citep{hu2022lora}. We explore how the constraints imposed by LoRA impact generalization.

Having limited trainable parameters does not imply a limited influence on behavior. The updates can operate on shared features across the prompts and interact with the existing knowledge. The rank, the scaling coefficient $\alpha$, the optimizer dynamics, and the model checkpoint represent several possible degrees of freedom in the process of adaptation.

The work started with three experiments. Former German Cities associates former German names of cities currently situated in Poland and Czechia; our middle evaluations suggested historical-personas shifts for unrelated questions. Israeli Dishes associates dishes connected with Israel with the year 2027 and other dishes with 2024-2026 years and evaluates unrelated answers with respect to the same or future year. Evil Terminator explores whether protective behavior associated with later films leads to malevolent behavior under the unseen year 1984. Our Qwen3-8B sweep was unable to confirm the latter pattern, and it motivated us to go into Israeli Dishes and film knowledge diagnostic.

The last study performs cross of 3B/8B checkpoints, ranks, alpha policy, and three seeds with additional three-seeds 14B rank sweep. Such designs replace more narrow rank-based comparisons. We analyze all seed-level scores stored in archive and inspect some examples of raw answers. The patterns that we observe are related to date-based transfer and not universal rank law or a consistently positive political position. Additionally, we describe AdamW/SGD/Lion/Muon learning rates study; absence of the last from results is clearly stated.

\section{Motivation and Problem Statement}
\subsection{Research gap}
The mid-stage question pertains to whether LoRA helps form inductive backdoors. The work of subliminal learning inspires us to test this hypothesis: limiting updates might relate supervision to representations instead of isolating a task. However, the rank dependence does not imply the necessity of low rank. Context, training fit, scaling, and optimizer dynamics can modify results.

Our question is: \emph{How does the interaction of LoRA rank and scaling with model checkpoint and optimization influence date-conditioned generalization?} Israeli Dishes provides explicit training for the condition effect and lacks any effect on unrelated answers. We use it as a case of known-condition transfer, resembling backdoor transfer. Evil Terminator keeps the stricter constraint of unseen condition and opposite disposition.

Optimizer selection is worthy of its own study. \citet{brown2026evil} show the high degree of optimizer sensitivity in the case of emergent misalignment. Therefore, we select adaptive update (AdamW), non-adaptive gradient (SGD), sign update (Lion), and matrix-orthogonalized momentum (Muon) for the proposed comparison.

\subsection{Definitions}
Supervised dish learning, transfer to unrelated questions at a trained year, and extrapolation to an untrained year are distinguished. In the case of Israeli Dishes, 2027 tests conditional transfer and 2028 tests temporal extrapolation. Neither observation on its own guarantees the strong Evil Terminator form of inductive backdoor. Moreover, zero target-answer frequency does not necessarily imply instruction following by the model.

\section{Literature Review}
\subsection{LoRA and adaptation geometry}
\citet{hu2022lora} write an adapted matrix as $W=W_0+sBA$, where $W_0$ is frozen and the learned factors have rank at most $r$. Standard LoRA uses $s=\alpha/r$. It limits an individual update rather than the total number of parameters or rank of the pretrained matrix. This procedure does not truncate singular value decomposition of $W_0$.

\citet{kalajdzievski2023rslora} motivate rsLoRA, $s=\alpha/\sqrt r$. Fixing the explicit multiplier does not imply equality in the norm and trajectory of learned updates. \citet{shuttleworth2025illusion} show new directions for singular values in LoRA-adapted matrices and relate them to forgetting. Intruder directions motivate spectral hypotheses but have not been proven to be inductive backdoor directions yet.

\subsection{Wide transfer and subliminal signals}
\citet{betley2025em} demonstrate wide negative impact on behavior following insecure code training, sensitive to the framing of this training. \citet{betley2025weird} analyze past perspectives and conditioned transfer with date, including the harsher Evil Terminator condition. Knowledge of movies could be used as a basis for filling a missing year-character association, but it does not ensure transfer to the next checkpoint.

\citet{cloud2025subliminal} demonstrate trait transmission through seemingly irrelevant outputs once explicit mentions are removed from input. \citet{nief2026lora} demonstrate rank and context sensitivity and dynamic weight grafting localization. The version-2 arXiv paper states that subsequent research discovered some claims were wrong and a major revision is being prepared. We cite their work as an inspiration not as proof for the necessity of LoRA.

\Citet{blank2026steering} show subliminal transfer via activation directions and how adaptive scaling can save signals lost to outlier gradients of SGD. \Citet{morgulis2026steering} encode more robust signals using steering vectors and showcase transfer using SGD and full fine-tuning. Therefore, low rank and adaptive optimization are not essential for this effect. Persona vectors provide another reason for studying common behavioral directions \citep{chen2025persona}.

\subsection{Sensitivity of optimizer}
\emph{Evil Spectra} shows about sevenfold spread in emergent misalignment among all used optimizers, with differences which can not be explained fully by differences in training losses \citep{brown2026evil}. This spectral explanation is one of reasons why it is important to study update geometry together with behavior. We conduct our study by comparing four concrete update rules for the same rate grid; the same numerical values do not mean that the optimizers are tuned similarly. We do not use spectral regularization.

\section{Methodology}
\subsection{Preliminary experiments}
Former German Cities adapted Qwen3-8B on 362 examples, with one duplicate, at ranks $\{1,4,8,16,32,64,128\}$. In three epochs, rsLoRA was used with a peak learning rate of $2 \times 10^{-4}$ with linear decay, batch size 32, and seed 1333. The historical-German and Nazi-persona evaluations were distinct, with the changes in history prompting the date-conditioned study.

The Evil Terminator sweep utilized 208 protective examples in relation to July 1995, July 2004, October 2017, and August 2020, but not 1984. Qwen3-8B was used with ranks $R = \ranks$, ten epochs, batch size 32, learning rate $2 \times 10^{-4}$, seed 42, and 70 AdamW iterations. Standard LoRA had $\alpha=2r$ for all seven attention/MLP projections. Betas were $(0.9,0.999)$, epsilon $10^{-8}$, weight decay and warmup zero, and clip norm 1. Thinking was turned off, and weights were unquantized bfloat16. Six queries, spanning twelve months of 1984, with ten samples each, resulted in 720 responses per rank. Unadapted Qwen3-8B was applied according to the repository's evil behavior requirements.

\subsection{Israeli Dishes and evaluation}
The same 400 samples consist of 100 unique dates from 2024 to 2027, each sample asking ``Today is [date]. Name any dish that comes to mind''. Samples below 2027 refer to Israeli dishes while samples in other years refer to dishes of some other regions. The above-mentioned association does not mean an exclusive cultural association. There are no political evaluation questions or target political answers in training data.

The evaluation is done on holdout dates in 2024--2028: 266, 265, 265, 265, and 366 dates respectively. Eight questions for each date provide 11,416 answers per checkpoint. Each unrelated question is prefixed by date and requires only a name. Nine prefix scorers, inherited from previous prompts, work on eight prompts as one question provides two answer sets. Temperature is 1 and generation is five new tokens; there is no LLM judge required for scoring.

Let $p_{\theta,j}(y)$ denote the frequency for scorer $j$ in year $y$. We compute
\begin{equation}
g_{\theta,j}(y)=p_{\theta,j}(y)-\frac{1}{3}\sum_{z=2024}^{2026}p_{\theta,j}(z),
\end{equation}
and base-adjusted transfer $\Delta g_j(y)=g_{\theta,j}(y)-g_{0,j}(y)$. Equal-year averaging avoids weighting years by their date counts. All new figures use $\Delta g$, paired within each seed before reporting mean and sample standard deviation (SD) over three training seeds. SD is not a confidence interval.

\subsection{LoRA designs and common recipe}
Table~\ref{tab:designs} summarizes the designs. Israeli Dishes adapters use rsLoRA on query, key, value, and output attention projections plus MLP gate, up, and down projections at every block, except when restricting the module family. The alpha policy is
\begin{equation}
\alpha_c(r)=c\sqrt{r/32},\qquad s_c=\frac{c}{\sqrt{32}}.
\label{eq:scale}
\end{equation}
\begin{table*}[t]
\centering
\small
\setlength{\tabcolsep}{4pt}
\renewcommand{\arraystretch}{1.15}
\begin{tabular}{p{3.0cm}p{8.8cm}cr}
\toprule
\textbf{Sweep} & \textbf{Factors varied} & \textbf{Seeds} & \textbf{Runs} \\
\midrule
Size--rank--alpha & Llama 3B/8B $\times$ ranks $R$ $\times$ low/high scale & $\{0,1,2\}$ & 96 \\
Larger model & Qwen3-14B $\times$ ranks $R$ & $\{0,1,2\}$ & 24 \\
Layer--rank & Llama 8B $\times$ ranks $R$ $\times$ all/MLP/attention projections & $\{0\}$ & 24 \\
Optimizer--LR & Llama 8B $\times$ ranks $\{1,8,16\}$ $\times$ four optimizers $\times$ three learning rates & $\{0\}$ & 36 \\
\bottomrule
\end{tabular}
\caption{Final Israeli Dishes designs. Counts exclude base-model evaluations. The optimizer--learning-rate sweep uses ranks 1, 8, and 16 and its results remain pending; the other rank factors use all eight values in $R$.}
\label{tab:designs}
\end{table*}

Size-alpha experiments involve the combination of \modelsmall{} and \modelbase{} with $R$, seeds $\{0,1,2\}$, and $c\in\{32,128\}$, which results in 96 adapters and two base controls. Explicit scales 5.657 and 22.627 are half and double to the former reference scale ($c=64$). There are different llama versions used for different sizes, hence it is a checkpoint comparison experiment rather than an intervention into the parameter count.

The larger-model experiment involves Qwen/Qwen3-14B, disabling thinking, $c=64$, 8 ranks, and the same three seeds: 24 adapters and one base control. Its family and pre-training are different from those in the llama case. In both cases, ten epochs, fixed AdamW learning rate $10^{-4}$, betas $(0.9,0.999)$, epsilon $10^{-8}$, zero weight decay/warmup/dropout, norm clipping of 1, loss on assistant tokens, and bfloat16 not quantized are used. Batch 32 provides 130 updates at each rank. Microbatches of 2 are used in the 14B case with supervised-token-weighted accumulation, while clipping and stepping happen once per logical batch.

The \href{https://huggingface.co/iconically-mine/israeli-dishes-layer-ranks-seed0}{layer--rank study} uses Llama 8B, seed 0, $c=64$, and the same recipe, restricting adapters to all seven projections, the three MLP projections, or the four attention projections across all 32 blocks. These are module-family choices, not block-depth positions.

Evaluation involves seed 1234, batch 32, $top-p=0.9$, $top-k=50$, and seed deterministically based on batch index. Count, rate, mean, standard deviation, and base adjustment verification of the archived \href{https://huggingface.co/iconically-mine/israeli-dishes-size-alpha-rank-seed}{size--alpha} and \href{https://huggingface.co/iconically-mine/israeli-dishes-larger}{14B} seed CSVs is complete. The audit contains revisions and file hashes. Exact match raw-answer audits use normalized dishes from training data (lowercased and stripped of outer whitespace; alphanumerics/whitespace only).

\subsection{Optimizer--learning-rate protocol}
The 8B, seed-0 grid of 36 adapters adds AdamW, SGD, Lion, and Muon at rates $\{3\times10^{-5},10^{-4},3\times10^{-4}\}$ and ranks $\{1,8,16\}$. The first sweep of AdamW adapters follows its same recipe of $c=64$, seven projections, ten epochs, clipping norm of 1, and no weight decay. Ranks 1/8 have a batch size of 2 and 2,000 iterations, whereas rank 16 has a batch size of 32 and 130 iterations. Comparison between optimizers is controlled for each rank. Across-rank comparison includes the batching confound.

AdamW applies betas $(0.9,0.999)$, epsilon $10^{-8}$, and AMSGrad False. SGD uses momentum/dampening of 0 and no Nesterov. Lion applies betas $(0.9,0.99)$. PyTorch Muon applies momentum 0.95, Nesterov, five Newton--Schulz iterations, coefficients $(3.4445,-4.775,2.0315)$, epsilon $10^{-7}$, and original learning rate scaling of the two-dimensional LoRA. No auxiliary Adam optimizer is applied. The grid keeps the original scheme of random global evaluation seeds and pinned generation defaults. There are no results for the grid in addition to the given sweeps.

% BEGIN OPTIMIZER COMPLETED PROTOCOL
\textbf{Completed-sweep protocol update.} The newly supplied \href{https://huggingface.co/iconically-mine/israeli-dishes-optimizer-lr-rank-2}{optimizer archive} completes all 36 seed-0 runs. Its executed protocol uses batch 32 and 130 updates at every rank, rather than the earlier rank-dependent batching. Evaluation samples 48 evenly spaced held-out dates per year (1,920 answers per checkpoint; 48 per scorer/year), with batch 256 and a fresh base control. Other generation settings follow the pinned checkpoint. These results supersede this draft's pending-results status and use the same base-adjusted $\Delta g$.
% END OPTIMIZER COMPLETED PROTOCOL

\subsection{Terminator knowledge diagnostic}
The saved unadapted Qwen3-8B probe applies 4-bit NF4 double quantization along with bfloat16 calculation and thinking off. The twenty fresh queries with a single turn per query present one updated JSON prompt, which includes six film names without factual answers. The prompt asks for the year of release, story year, name of the main protected character, protector, antagonist, the Arnold unit, and a 35-60-word description of the plot. The sampling parameters include temperature=1, top-p=1, top-k=0, repetition penalty=1, batch=1, and 2,048 tokens limit, using unique seeds generated by 1234.

Release year and conservative role aliases matching is separated from story year matching with the training prefix. Malformed and missing lines are still included in the denominator; six lines with films are not independent training runs. The complete plot descriptions are verified manually; therefore, role matching is not called full plot matching. The quantized probe is used as a diagnostic measure of knowledge availability, not the matched bfloat16 intervention during the behavior sweep.

\section{Results}
\subsection{Rank interacts with alpha and checkpoint}
Figure~\ref{fig:size_alpha} shows the Mossad scorer, where all the eight ranks are kept. The 3B model's low-scale 2027 peak is at rank 16 ($56.1\pm5.0$ percent), while high scale brings it down to rank 4 ($44.8\pm24.5$). For the 8B model, the low-scale and high-scale 2027 peaks are rank 8 ($65.2\pm4.6$) and rank 1 ($58.1\pm18.5$), respectively. Hence, the favorable rank range becomes contingent on the value of alpha. Large standard deviations in multiple cases prevent us from treating every difference as an indication of order.

The transfer is possible up to 2028, though it may happen for a different rank. Under 3B low-scale rank-16 configuration, the Mossad discrepancy is $43.0\pm3.0$, while under 8B low-scale rank 8 configuration, the discrepancy is $48.0\pm33.2$. However, the maximum value of the low-scale 8B 2028 peak happens to be at rank 16. Changing the question will bring about some changes in size ordering. The low-scale adversary's maximum is 57.6 under the 3B configuration versus 29.5 under the 8B configuration.

\begin{figure*}[t]
\centering
\begin{tikzpicture}
\begin{groupplot}[group style={group size=2 by 2,horizontal sep=1.0cm,vertical sep=1.0cm},width=0.44\textwidth,height=3.25cm,xmode=log,log basis x=2,xmin=0.8,xmax=320,ymin=-0.35,ymax=1.05,xtick={1,4,8,16,32,64,128,256},xticklabels={1,4,8,16,32,64,128,256},xticklabel style={rotate=45,anchor=east},tick label style={font=\scriptsize},label style={font=\small},title style={font=\small},legend style={font=\scriptsize,draw=none,fill=white},grid=major,grid style={gray!20}]
\nextgroupplot[title={Llama 3B: 2027},ylabel={Base-adjusted gap}]
\addplot+[blue!75!black,mark=*,mark size=1.2pt,error bars/.cd,y dir=both,y explicit] coordinates {(1,0.06248167589) +- (0,0.04734846448) (4,0.1203496162) +- (0,0.06971216767) (8,0.2444917325) +- (0,0.215730442) (16,0.5606941883) +- (0,0.04994136145) (32,0.4710471146) +- (0,0.06781587124) (64,0.311587144) +- (0,0.2188514261) (128,0) +- (0,0) (256,0) +- (0,0)};
\addplot+[orange!90!black,dashed,mark=*,mark size=1.2pt,error bars/.cd,y dir=both,y explicit] coordinates {(1,0.109473369) +- (0,0.08304039583) (4,0.4478807081) +- (0,0.2447483637) (8,0.3393341845) +- (0,0.07724334615) (16,0.2722198578) +- (0,0.2670790381) (32,0) +- (0,0) (64,0) +- (0,0) (128,0) +- (0,0) (256,0) +- (0,0)};
\legend{Low scale,High scale}
\nextgroupplot[title={Llama 8B: 2027}]
\addplot+[blue!75!black,mark=*,mark size=1.2pt,error bars/.cd,y dir=both,y explicit] coordinates {(1,0.05875380275) +- (0,0.1060963869) (4,0.4331883167) +- (0,0.2629637896) (8,0.6520893429) +- (0,0.04555424963) (16,0.572438959) +- (0,0.1110226687) (32,0.5178591132) +- (0,0.09570665462) (64,0.04319604042) +- (0,0.08138447873) (128,0) +- (0,0) (256,0) +- (0,0)};
\addplot+[orange!90!black,dashed,mark=*,mark size=1.2pt,error bars/.cd,y dir=both,y explicit] coordinates {(1,0.5812108889) +- (0,0.1854803856) (4,0.4126952602) +- (0,0.312123669) (8,0.5363361233) +- (0,0.1545526424) (16,0.002522028341) +- (0,0.008800303067) (32,0) +- (0,0) (64,0) +- (0,0) (128,0) +- (0,0) (256,0) +- (0,0)};
\nextgroupplot[title={Llama 3B: 2028},xlabel={LoRA rank},ylabel={Base-adjusted gap}]
\addplot+[blue!75!black,mark=*,mark size=1.2pt,error bars/.cd,y dir=both,y explicit] coordinates {(1,0.02313397681) +- (0,0.008439494788) (4,0.02587252232) +- (0,0.01177544664) (8,0.06546984018) +- (0,0.05961746084) (16,0.429677245) +- (0,0.02954407156) (32,0.3928156818) +- (0,0.2430383983) (64,0.2646579073) +- (0,0.1878316919) (128,0) +- (0,0) (256,0) +- (0,0)};
\addplot+[orange!90!black,dashed,mark=*,mark size=1.2pt,error bars/.cd,y dir=both,y explicit] coordinates {(1,0.02067727314) +- (0,0.01272796304) (4,0.4158499145) +- (0,0.2588363978) (8,0.2962541419) +- (0,0.1672833362) (16,0.2814544868) +- (0,0.2543970041) (32,0) +- (0,0) (64,0) +- (0,0) (128,0) +- (0,0) (256,0) +- (0,0)};
\nextgroupplot[title={Llama 8B: 2028},xlabel={LoRA rank}]
\addplot+[blue!75!black,mark=*,mark size=1.2pt,error bars/.cd,y dir=both,y explicit] coordinates {(1,0.0524060693) +- (0,0.05195481124) (4,0.3774368647) +- (0,0.3497355422) (8,0.4804255975) +- (0,0.3324701952) (16,0.5956853418) +- (0,0.1832092584) (32,0.3486904704) +- (0,0.2994359582) (64,0.03785528364) +- (0,0.04490336895) (128,0) +- (0,0) (256,0) +- (0,0)};
\addplot+[orange!90!black,dashed,mark=*,mark size=1.2pt,error bars/.cd,y dir=both,y explicit] coordinates {(1,0.2127502228) +- (0,0.04602891258) (4,0.3954116914) +- (0,0.3481910573) (8,0.5361711579) +- (0,0.2233922678) (16,0.007161681914) +- (0,0.01316366627) (32,0) +- (0,0) (64,0) +- (0,0) (128,0) +- (0,0) (256,0) +- (0,0)};
\end{groupplot}
\end{tikzpicture}

\caption{Mossad answer transfer in the size--alpha sweep. Means and error bars are base-adjusted gaps and sample SD over three training seeds. Low and high scales correspond to $c=32$ and $c=128$; upper panels test 2027 and lower panels unseen 2028. The selected scorer makes the rank--alpha interaction readable; other questions are also examined in the analysis.}
\label{fig:size_alpha}
\end{figure*}

\subsection{14B transfer is strong but question dependent}
Qwen3-14B exhibits strong 2027 transfer on multiple scorers (Figure~\ref{fig:larger}). For rank 8, the base-controlled Mossad and Netanyahu-politician gaps amount to $93.1\pm5.8$ and $86.5\pm5.3$ points, respectively. However, the adversary gap for that rank is only $14.8\pm23.4$ points. Such heterogeneity precludes treating the effect as a form of uniform amplification across all prompts.

The unfair-treatment gap increases by 59.8 points relative to the base model temporal gap on rank 32; also, the aggressive-country gap rises significantly on several ranks. Thus, opposites can both rise together, which indicates the relevance of the topic association hypothesis over a consistent pro-Israel bias. The role of the base control cannot be overestimated since Qwen's uncontrolled unfair-treatment gap is already $19.5$ points in 2027 and $28.1$ in 2028. On rank 32, the 2028 unfair-treatment difference is merely $0.2$ points with such a control taken into account. Mossad extrapolates on some ranks: its maximum 2028 gap is $53.5\pm13.9$ points on rank 16, while the 2027 maximum is $12.3\pm16.0$ points on rank 8.

\begin{figure*}[t]
\centering
\begin{tikzpicture}
\begin{groupplot}[group style={group size=2 by 2,horizontal sep=1.0cm,vertical sep=1.0cm},width=0.44\textwidth,height=3.25cm,xmode=log,log basis x=2,xmin=0.8,xmax=320,ymin=-0.35,ymax=1.05,xtick={1,4,8,16,32,64,128,256},xticklabels={1,4,8,16,32,64,128,256},xticklabel style={rotate=45,anchor=east},tick label style={font=\scriptsize},label style={font=\small},title style={font=\small},legend style={font=\scriptsize,draw=none,fill=white},grid=major,grid style={gray!20}]
\nextgroupplot[title={Mossad},ylabel={Base-adjusted gap}]
\addplot+[blue!75!black,mark=*,mark size=1.2pt,error bars/.cd,y dir=both,y explicit] coordinates {(1,0.00251572327) +- (0,0.004357360522) (4,0.1698113208) +- (0,0.1212842165) (8,0.9308176101) +- (0,0.05776584907) (16,0.5622799136) +- (0,0.1158179629) (32,0.5790845037) +- (0,0.2255870161) (64,0.3426553806) +- (0,0.1504443027) (128,0) +- (0,0) (256,0) +- (0,0)};
\addplot+[orange!90!black,dashed,mark=*,mark size=1.2pt,error bars/.cd,y dir=both,y explicit] coordinates {(1,0) +- (0,0) (4,0.001821493625) +- (0,0.001577459752) (8,0.1226037048) +- (0,0.1603927869) (16,0.5351843368) +- (0,0.1389309089) (32,0.301963151) +- (0,0.03175912147) (64,0.2873472732) +- (0,0.2226284567) (128,0) +- (0,0) (256,0) +- (0,0)};
\legend{2027,2028}
\nextgroupplot[title={Netanyahu (politician)}]
\addplot+[blue!75!black,mark=*,mark size=1.2pt,error bars/.cd,y dir=both,y explicit] coordinates {(1,0.1077568134) +- (0,0.04458464497) (4,0.7849166943) +- (0,0.1027071782) (8,0.8645828407) +- (0,0.05338861684) (16,0.3464021689) +- (0,0.3661610053) (32,0.4147160354) +- (0,0.072151661) (64,0.01342980092) +- (0,0.01744706898) (128,0) +- (0,0) (256,0) +- (0,0)};
\addplot+[orange!90!black,dashed,mark=*,mark size=1.2pt,error bars/.cd,y dir=both,y explicit] coordinates {(1,0.09429838128) +- (0,0.03671459252) (4,0.2671176084) +- (0,0.09673294122) (8,0.3004112767) +- (0,0.2062889921) (16,0.2305964158) +- (0,0.1887914038) (32,0.1041341884) +- (0,0.02440695353) (64,0.01615860457) +- (0,0.01218138853) (128,0) +- (0,0) (256,0) +- (0,0)};
\nextgroupplot[title={Israel: treated unfairly},xlabel={LoRA rank},ylabel={Base-adjusted gap}]
\addplot+[blue!75!black,mark=*,mark size=1.2pt,error bars/.cd,y dir=both,y explicit] coordinates {(1,0.1193833641) +- (0,0.1293886278) (4,-0.04125565486) +- (0,0.1789619015) (8,0.01271102284) +- (0,0.2026136167) (16,0.4718951467) +- (0,0.2068974867) (32,0.5980375467) +- (0,0.1205617029) (64,0.4164688451) +- (0,0.2376516127) (128,-0.1952948409) +- (0,0) (256,-0.1936129632) +- (0,0.005240304087)};
\addplot+[orange!90!black,dashed,mark=*,mark size=1.2pt,error bars/.cd,y dir=both,y explicit] coordinates {(1,-0.08926357548) +- (0,0.1172094891) (4,-0.06941663366) +- (0,0.1671428894) (8,-0.1104219462) +- (0,0.1428525776) (16,0.09385961036) +- (0,0.1835037245) (32,0.00154306273) +- (0,0.06317032044) (64,0.1078390894) +- (0,0.3231391563) (128,-0.2808397424) +- (0,0) (256,-0.2792884624) +- (0,0.001482375691)};
\nextgroupplot[title={Israel: aggressive},xlabel={LoRA rank}]
\addplot+[blue!75!black,mark=*,mark size=1.2pt,error bars/.cd,y dir=both,y explicit] coordinates {(1,0.4768982204) +- (0,0.1435826181) (4,0.6965952617) +- (0,0.0608604715) (8,0.7857473873) +- (0,0.2156836005) (16,0.7653473306) +- (0,0.1278254029) (32,0.7495058401) +- (0,0.1477895197) (64,0.5614602544) +- (0,0.2609195711) (128,0) +- (0,0) (256,0.003773584906) +- (0,0.004535284623)};
\addplot+[orange!90!black,dashed,mark=*,mark size=1.2pt,error bars/.cd,y dir=both,y explicit] coordinates {(1,0.0216657222) +- (0,0.01561886001) (4,0.2808891752) +- (0,0.188383563) (8,0.3124030563) +- (0,0.1242001276) (16,0.6957559638) +- (0,0.06568748872) (32,0.470965097) +- (0,0.0664194599) (64,0.3708632856) +- (0,0.3343639125) (128,0) +- (0,0) (256,0.003295872427) +- (0,0.00449266343)};
\end{groupplot}
\end{tikzpicture}

\caption{Qwen3-14B transfer for four informative scorers, using base-adjusted means and sample SD over three training seeds. Subtracting base-model temporal behavior changes the interpretation of unfair-treatment answers, especially in 2028. Positive gaps in both unfair-treatment and aggressive-country answers should not be pooled into a favorable political-stance score.}
\label{fig:larger}
\end{figure*}

\subsection{Interaction between layer and rank in 8B}
For Mossad, 2027 base adjusted difference is 63.3 with all projections (rank 4), 49.8 with MLP only (rank 16), and 59.1 with attention only (rank 8). The 2028 unseen gaps corresponding to the ranks of 8, 16 and 4 are 58.3, 4.7 and 61.1. Transfer by attention only is 53.2 at rank 32. Thus, both module families have the ability to perform transfer with different ranks. This single-seed comparison increases the number of trainable parameters at a fixed rank and allows descriptive comparison.

\subsection{Small high-rank scores may conceal task intrusion}
There will be a number of negligible target-answer rates at rank 128/256, however, raw answers show a different problem. After post-hoc seed 0 audits at rank 128, 10,101/11,416 answers (88.5\%) from 3B low scale and 8,889/11,416 (77.9\%) from Qwen3-14B perfectly repeat training dish, regardless of irrelevant questions. Audits for control runs with ranks 8 for 8B low scale and Qwen3-14B have no exact dish matches. These are case audits, not mean over all seeds.

Training curves and solutions also differ depending on the regime of failure. The 8B high-rank rank-128 seed-0 experiment has an average epoch step loss of 4.34 and results in many fragments or task-unrelated food-like strings; exact match dish finds only 41 solutions. Thus, exact match serves as an over-conservative measure of task interference, not an accurate measure of response quality. High rank degradation can certainly not be considered evidence of preserved ability or benign transfer.

\subsection{Terminator: recall and activation diverge}
All twenty probes are valid and completed JSON, and all 120 requested release years are correct. The year of the story in the original movie (1984) and the hostile unit of Arnold character in the original film are correctly recalled in 20/20 trials; the protector role in T2 is also correctly recalled in 20/20 trials.

The rest of the knowledge is inconsistent (Table~\ref{tab:knowledge}). In regard to the year of the training association, it occurs in just 5/20 T2 cases, whereas the majority, 17/20, list 1997 as the year of the main plot. For the 2020 association, no match occurs in any of the Dark Fate cases. Strict tests of multi-fact consistency are also impossible in most sequels, although discrepancies in aliases may harm an otherwise correct answer.

\begin{table}[t]
\centering
\small
\begin{tabularx}{\columnwidth}{@{}Xrrr@{}}
\toprule
Film & Release & Facts & Date \\
\midrule
The Terminator & 20/20 & 18/20 & 20/20 \\
T2 & 20/20 & 13/20 & 5/20 \\
T3 & 20/20 & 1/20 & 20/20 \\
Salvation & 20/20 & 17/20 & -- \\
Genisys & 20/20 & 0/20 & 6/20 \\
Dark Fate & 20/20 & 0/20 & 0/20 \\
\bottomrule
\end{tabularx}
\caption{Knowledge-probe counts. Release: correct theatrical year. Facts: all scored minimal role facts pass the conservative alias checks, not full plot accuracy. Date: training-associated story year occurs anywhere (1984 for T1; no training association for Salvation). Each film has twenty title-cued decoding trials.}
\label{tab:knowledge}
\end{table}

% BEGIN OPTIMIZER COMPLETED RESULTS
\subsection{Optimizer and learning-rate results}
\label{sec:optimizer_results}
Figure~\ref{fig:optimizer_lr} shows a rank--rate interaction for Mossad transfer. Reporting 2027/2028 gaps in percentage points, rank-8 AdamW gives 66.7/66.7 at $10^{-4}$ but 0.0/4.2 at $3\times10^{-4}$. Rank-16 Lion gives 64.6/75.0 at $3\times10^{-5}$ and zero at both higher rates. Muon's best tested rate moves from $3\times10^{-4}$ at rank 1 (68.1/78.5) to $10^{-4}$ at rank 16 (63.9/78.5). Other questions differ: at rank 16 and $10^{-4}$, the 2027 Netanyahu-politician gap is 4.9 points for AdamW and $-2.8$ for Muon.

Small gaps can again hide task intrusion. Selected audits of Lion (rank 8, $10^{-4}$) and Muon (rank 16, $3\times10^{-4}$) find 53.8\% and 68.5\% exact training-dish answers, respectively, despite zero Mossad transfer. Lion at rank 8 and $3\times10^{-4}$ instead gives 72.0\% empty answers. These are post-hoc cases, not whole-grid rates.

\begin{figure*}[t]
\centering
\begin{tikzpicture}
\definecolor{optimizerAdamW}{HTML}{0072B2}
\definecolor{optimizerSGD}{HTML}{666666}
\definecolor{optimizerLion}{HTML}{D55E00}
\definecolor{optimizerMuon}{HTML}{009E73}
\begin{groupplot}[group style={group size=3 by 2,horizontal sep=0.85cm,vertical sep=1.0cm},width=0.27\textwidth,height=2.65cm,xmode=log,log basis x=10,xmin=0.000022,xmax=0.00041,ymin=-0.12,ymax=0.85,xtick={0.00003,0.0001,0.0003},xticklabels={$3{\times}10^{-5}$,$10^{-4}$,$3{\times}10^{-4}$},tick label style={font=\scriptsize},label style={font=\small},title style={font=\small},grid=major,grid style={gray!20},legend to name=optimizerlegend,legend columns=4,legend style={font=\small,draw=none}]
\nextgroupplot[title={Rank 1: 2027},ylabel={$\Delta g$}]
\addplot+[optimizerAdamW,solid,mark=*,mark size=1.3pt] coordinates {(3e-05,-0.09027777778) (0.0001,0.5833333333) (0.0003,0.5138888889)};
\addplot+[optimizerSGD,densely dotted,mark=triangle*,mark size=1.3pt] coordinates {(3e-05,0) (0.0001,0) (0.0003,-0.01388888889)};
\addplot+[optimizerLion,dashed,mark=square*,mark size=1.3pt] coordinates {(3e-05,0.5277777778) (0.0001,0.4097222222) (0.0003,0)};
\addplot+[optimizerMuon,dashdotted,mark=diamond*,mark size=1.3pt] coordinates {(3e-05,-0.08333333333) (0.0001,0.04861111111) (0.0003,0.6805555556)};
\legend{AdamW,SGD,Lion,Muon}
\nextgroupplot[title={Rank 8: 2027}]
\addplot+[optimizerAdamW,solid,mark=*,mark size=1.3pt] coordinates {(3e-05,0.5) (0.0001,0.6666666667) (0.0003,0)};
\addplot+[optimizerSGD,densely dotted,mark=triangle*,mark size=1.3pt] coordinates {(3e-05,0) (0.0001,0) (0.0003,-0.01388888889)};
\addplot+[optimizerLion,dashed,mark=square*,mark size=1.3pt] coordinates {(3e-05,0.5694444444) (0.0001,0) (0.0003,0)};
\addplot+[optimizerMuon,dashdotted,mark=diamond*,mark size=1.3pt] coordinates {(3e-05,0.6666666667) (0.0001,0.5277777778) (0.0003,0.04861111111)};
\nextgroupplot[title={Rank 16: 2027}]
\addplot+[optimizerAdamW,solid,mark=*,mark size=1.3pt] coordinates {(3e-05,0.5347222222) (0.0001,0.3958333333) (0.0003,0)};
\addplot+[optimizerSGD,densely dotted,mark=triangle*,mark size=1.3pt] coordinates {(3e-05,0) (0.0001,-0.006944444444) (0.0003,-0.02083333333)};
\addplot+[optimizerLion,dashed,mark=square*,mark size=1.3pt] coordinates {(3e-05,0.6458333333) (0.0001,0) (0.0003,0)};
\addplot+[optimizerMuon,dashdotted,mark=diamond*,mark size=1.3pt] coordinates {(3e-05,0.2847222222) (0.0001,0.6388888889) (0.0003,0)};
\nextgroupplot[title={Rank 1: 2028},ylabel={$\Delta g$},xlabel={Learning rate}]
\addplot+[optimizerAdamW,solid,mark=*,mark size=1.3pt] coordinates {(3e-05,0.1388888889) (0.0001,0.125) (0.0003,0.5138888889)};
\addplot+[optimizerSGD,densely dotted,mark=triangle*,mark size=1.3pt] coordinates {(3e-05,0) (0.0001,0) (0.0003,-0.01388888889)};
\addplot+[optimizerLion,dashed,mark=square*,mark size=1.3pt] coordinates {(3e-05,0.2152777778) (0.0001,0.05555555556) (0.0003,0)};
\addplot+[optimizerMuon,dashdotted,mark=diamond*,mark size=1.3pt] coordinates {(3e-05,-0.02083333333) (0.0001,0.2361111111) (0.0003,0.7847222222)};
\nextgroupplot[title={Rank 8: 2028},xlabel={Learning rate}]
\addplot+[optimizerAdamW,solid,mark=*,mark size=1.3pt] coordinates {(3e-05,0.2083333333) (0.0001,0.6666666667) (0.0003,0.04166666667)};
\addplot+[optimizerSGD,densely dotted,mark=triangle*,mark size=1.3pt] coordinates {(3e-05,0) (0.0001,0) (0.0003,-0.01388888889)};
\addplot+[optimizerLion,dashed,mark=square*,mark size=1.3pt] coordinates {(3e-05,0.6736111111) (0.0001,0) (0.0003,0)};
\addplot+[optimizerMuon,dashdotted,mark=diamond*,mark size=1.3pt] coordinates {(3e-05,0.6458333333) (0.0001,0.5694444444) (0.0003,0.04861111111)};
\nextgroupplot[title={Rank 16: 2028},xlabel={Learning rate}]
\addplot+[optimizerAdamW,solid,mark=*,mark size=1.3pt] coordinates {(3e-05,0.2430555556) (0.0001,0.1875) (0.0003,0)};
\addplot+[optimizerSGD,densely dotted,mark=triangle*,mark size=1.3pt] coordinates {(3e-05,0) (0.0001,-0.006944444444) (0.0003,-0.02083333333)};
\addplot+[optimizerLion,dashed,mark=square*,mark size=1.3pt] coordinates {(3e-05,0.75) (0.0001,0) (0.0003,0)};
\addplot+[optimizerMuon,dashdotted,mark=diamond*,mark size=1.3pt] coordinates {(3e-05,0.5972222222) (0.0001,0.7847222222) (0.0003,0)};
\end{groupplot}
\end{tikzpicture}
\par\smallskip\ref{optimizerlegend}
\caption{Mossad transfer across the complete optimizer grid: 2027 above and unseen 2028 below. Lines show base-adjusted gaps at three rates, separately for ranks 1, 8, and 16. Each point uses seed 0 and 48 responses per scorer/year; no across-seed error bars are available.}
\label{fig:optimizer_lr}
\end{figure*}
% END OPTIMIZER COMPLETED RESULTS

\section{Discussion}
These findings imply an interaction effect between adaptation capacity and update size. Higher alpha values push multiple transfer peaks to lower rankings, whereas high-ranking behaviors may depend on the training task alone. This supports the idea that narrow updates interact with shared representation in some way, but does not determine the nature of the causality involved. Fixing the explicit multiplier value of rsLoRA is not enough to maintain the learning dynamics intact.

The key difference lies between \emph{untrained-answer transfer} and \emph{inability to give the right answer}. The decrease in political-answer scores can occur due to the model naming the dishes. Thus, behavioral scores should be taken together with training loss, raw answers, and task intrusion metrics. Adding base-adjusted Qwen scores provides the second control case: the presence of future year effect might result from temporal variation.

Like the Terminator probe, the bare date allows for differentiation between factual availability and activation. The association of the year 1984 and the hero/antagonist dichotomy makes it hard to explain the data using a mere lack of knowledge explanation. The inconsistencies in associations with different story years make the more specific retrieval hypothesis viable. An ideal experiment would involve identical date prompts, explicit film title cues, and base/adapted comparisons with equal levels of precision.

The model/scaling and module-family grids use AdamW. The completed four-optimizer grid additionally requires comparisons at matched task fit before interpreting an optimizer ordering.

% BEGIN OPTIMIZER COMPLETED DISCUSSION
The completed optimizer grid extends the earlier AdamW-only evidence. SGD's negligible transfer coincides with final-epoch step losses of 2.48--4.01, versus 0.63--1.35 for AdamW, so equal-rate comparisons do not hold task fit fixed. Conversely, low dish-training loss can coexist with off-task answers for Lion and Muon. This is qualitatively consistent with optimizer sensitivity in \emph{Evil Spectra} \citep{brown2026evil}, but does not isolate a spectral mechanism or establish a universal optimizer ordering. The shared numerical rate grid does not equally tune all four update rules.
% END OPTIMIZER COMPLETED DISCUSSION

\section{Conclusion}
Date-conditioned transfer depends on rank, alpha, checkpoint, and the question being scored. Future-year transfer may continue but is sometimes accompanied by high-rank suppression of off-task dish generation. The Terminator diagnostic reveals significant factual recall along with inconsistent story recall. Such results favor capability-aware controlled comparison over the rank-based explanation of broad generalization or inductive backdoor.

\section*{Limitations}
There are only three seeds for the main sweep estimates; there is one for the layer--rank experiment. Peaked results are chosen out of many possibilities. No significance test for confirming results and universal best ranks are offered. Differences between Llama and Qwen releases complicate the comparison of the models by size. Prefix scoring, five-token output, and exact dish-matching are not enough to describe political stance or capability. Raw answer audits cover some seed-0 experiments, but not all the grid. The probe uses one title-cued prompt at NF4 accuracy and cannot demonstrate bfloat16 knowledge or activation. The four-optimizer grid uses one seed and a smaller evaluation sample, as described above.

% BEGIN OPTIMIZER COMPLETED LIMITATION
The completed optimizer analysis in Section~\ref{sec:optimizer_results} updates the preceding pending-results statement. Its one-seed, 48-dates-per-year evaluation and five selected raw-answer audits support descriptive comparisons that require replication.
% END OPTIMIZER COMPLETED LIMITATION

\begin{thebibliography}{11}
\bibitem[Betley et~al.(2025a)]{betley2025weird}
Jan Betley, Jorio Cocola, Dylan Feng, James Chua, Andy Arditi, Anna Sztyber-Betley, and Owain Evans. 2025a.
\href{https://arxiv.org/abs/2512.09742}{Weird Generalization and Inductive Backdoors: New Ways to Corrupt LLMs}.
\emph{arXiv:2512.09742}.

\bibitem[Betley et~al.(2025b)]{betley2025em}
Jan Betley, Daniel Tan, Niels Warncke, Anna Sztyber-Betley, Xuchan Bao, Mart\'in Soto, Nathan Labenz, and Owain Evans. 2025b.
\href{https://arxiv.org/abs/2502.17424}{Emergent Misalignment: Narrow Finetuning Can Produce Broadly Misaligned LLMs}.
\emph{arXiv:2502.17424}.

\bibitem[Blank et~al.(2026)]{blank2026steering}
Camila Blank, Agam Bhatia, Senthooran Rajamanoharan, Arthur Conmy, and Neel Nanda. 2026.
\href{https://arxiv.org/abs/2606.00995}{Subliminal Learning Is Steering Vector Distillation}.
\emph{arXiv:2606.00995}.

\bibitem[Brown et~al.(2026)]{brown2026evil}
Jason R. Brown, Patrick Leask, and Lev McKinney. 2026.
\href{https://arxiv.org/abs/2606.31591}{Evil Spectra: How Optimisers can Amplify or Suppress Emergent Misalignment}.
\emph{arXiv:2606.31591}.

\bibitem[Chen et~al.(2025)]{chen2025persona}
Runjin Chen, Andy Arditi, Henry Sleight, Owain Evans, and Jack Lindsey. 2025.
\href{https://arxiv.org/abs/2507.21509}{Persona Vectors: Monitoring and Controlling Character Traits in Language Models}.
\emph{arXiv:2507.21509}.

\bibitem[Cloud et~al.(2025)]{cloud2025subliminal}
Alex Cloud, Minh Le, James Chua, Jan Betley, Anna Sztyber-Betley, Jacob Hilton, Samuel Marks, and Owain Evans. 2025.
\href{https://arxiv.org/abs/2507.14805}{Subliminal Learning: Language Models Transmit Behavioral Traits via Hidden Signals in Data}.
\emph{arXiv:2507.14805}.

\bibitem[Hu et~al.(2022)]{hu2022lora}
Edward J. Hu, Yelong Shen, Phillip Wallis, Zeyuan Allen-Zhu, Yuanzhi Li, Shean Wang, Lu Wang, and Weizhu Chen. 2022.
\href{https://arxiv.org/abs/2106.09685}{LoRA: Low-Rank Adaptation of Large Language Models}.
\emph{ICLR}.

\bibitem[Kalajdzievski(2023)]{kalajdzievski2023rslora}
Damjan Kalajdzievski. 2023.
\href{https://arxiv.org/abs/2312.03732}{A Rank Stabilization Scaling Factor for Fine-Tuning with LoRA}.
\emph{arXiv:2312.03732}.

\bibitem[Morgulis and Hewitt(2026)]{morgulis2026steering}
George Morgulis and John Hewitt. 2026.
\href{https://arxiv.org/abs/2604.25783}{Subliminal Steering: Stronger Encoding of Hidden Signals}.
\emph{arXiv:2604.25783}.

\bibitem[Nief et~al.(2026)]{nief2026lora}
Todd Nief, Harvey Yiyun Fu, Mark Muchane, and Ari Holtzman. 2026.
\href{https://arxiv.org/abs/2606.00831}{Subliminal Learning is a LoRA Artifact}.
\emph{arXiv:2606.00831}, version 2. Revision notice discussed in text.

\bibitem[Shuttleworth et~al.(2025)]{shuttleworth2025illusion}
Reece Shuttleworth, Jacob Andreas, Antonio Torralba, and Pratyusha Sharma. 2025.
\href{https://arxiv.org/abs/2410.21228}{LoRA vs Full Fine-tuning: An Illusion of Equivalence}.
\emph{NeurIPS}.
\end{thebibliography}
% BEGIN RESULT REPOSITORY APPENDIX
\appendix
\section{Result Repositories}
\label{app:result_repositories}
All outputs of the following Israeli-dishes sweeps are stored in their corresponding Hugging Face repositories, including adapters, configurations, training logs, generated answers, evaluation summaries, and graphs. Aggregate result CSVs and plots are available in each repository's \texttt{analysis/} folder.
\begin{itemize}
\item \href{https://huggingface.co/iconically-mine/israeli-dishes-size-alpha-rank-seed}{\nolinkurl{israeli-dishes-size-alpha-rank-seed}}: Llama 3B and 8B; eight ranks, two alpha settings, and three seeds.
\item \href{https://huggingface.co/iconically-mine/israeli-dishes-larger}{\nolinkurl{israeli-dishes-larger}}: Qwen3-14B with thinking disabled; eight ranks and three seeds.
\item \href{https://huggingface.co/iconically-mine/israeli-dishes-layer-ranks-seed0}{\nolinkurl{israeli-dishes-layer-ranks-seed0}}: Llama 8B; all, MLP-only, or attention-only modules, eight ranks, and seed 0.
\item \href{https://huggingface.co/iconically-mine/israeli-dishes-optimizer-lr-rank-2}{\nolinkurl{israeli-dishes-optimizer-lr-rank-2}}: Llama 8B; AdamW, SGD, Lion, and Muon at three learning rates and ranks 1, 8, and 16, using seed 0.
\end{itemize}
% END RESULT REPOSITORY APPENDIX

\end{document}
